% 第3章3.4.1下界平移：在三点算例中把 s->t 的下界改为1。
% 原网络：s->t 的界[1,4]、费用4；平移 x=l+y 后 s 的剩余供给 3-1=2，t 的剩余需求 -3+1=-2，
% s->t 容量 4-1=3，目标常数 4*1=4。最优 y=(2,2,0)，费用4，加常数得原费用8（verify.py）。
\begin{tikzpicture}[>=Stealth,line width=.8pt,font=\fontsize{8}{10}\selectfont,
  v/.style={circle,draw=FigureStroke,fill=FigureBlueFill,line width=.8pt,minimum size=6.5mm,inner sep=0pt,text=FigureInk},
  e/.style={->,draw=FigureStroke,line width=.8pt,shorten >=.5pt},
  lab/.style={text=FigureInk,inner sep=1.5pt},
  note/.style={text=FigureMuted,font=\FigureNoteFont\fontsize{7.5}{9}\selectfont}]
  % (a)
  \begin{scope}
  \node[v] (s) at (0,0) {$s$};
  \node[v] (a) at (1.8,1.4) {$a$};
  \node[v] (t) at (3.6,0) {$t$};
  \draw[e] (s)--node[lab,above left=-1pt]{$[0,3],1$} (a);
  \draw[e] (a)--node[lab,above right=-1pt]{$[0,2],1$} (t);
  \draw[e,draw=FigureYellow,line width=1.2pt] (s)--node[lab,below]{$[1,4],4$} (t);
  \node[text=FigureBlue,below=1pt] at (s.south) {$b_s=3$};
  \node[text=FigureBlue,below=1pt] at (t.south) {$b_t=-3$};
  \node at (1.8,-1.25) {(a) 原网络：弧旁为$[\ell,u],c$};
  \end{scope}
  \draw[->,FigureMuted,line width=.8pt] (4.35,.5)--node[above,note]{$x=\ell+y$} (5.45,.5);
  % (b)
  \begin{scope}[xshift=6.2cm]
  \node[v] (s) at (0,0) {$s$};
  \node[v] (a) at (1.8,1.4) {$a$};
  \node[v] (t) at (3.6,0) {$t$};
  \draw[e] (s)--node[lab,above left=-1pt]{$[0,3],1$} (a);
  \draw[e] (a)--node[lab,above right=-1pt]{$[0,2],1$} (t);
  \draw[e,draw=FigureYellow,line width=1.2pt] (s)--node[lab,below]{$[0,3],4$} (t);
  \node[text=FigureGreen,below=1pt] at (s.south) {$b'_s=2$};
  \node[text=FigureGreen,below=1pt] at (t.south) {$b'_t=-2$};
  \node[note,align=center] at (1.8,2.25) {目标另加常数$c_{s,t}\ell_{s,t}=4$};
  \node at (1.8,-1.25) {(b) 平移后：下界为$0$};
  \end{scope}
\end{tikzpicture}
